\chapter{The Statistical Mechanics of Vortices: Deriving Thermodynamics}

\author{James Tyndall}

%----------------------------------------------------------------------
\section{Introduction: From Single Atoms to Collective Behaviour}
%----------------------------------------------------------------------

The world we experience is not one of isolated atoms but of vast statistical ensembles.  SDT provides a direct, mechanical foundation: the laws of thermodynamics are the \textbf{emergent, macroscopic consequences of the interactions between countless displacement vortices.}


%----------------------------------------------------------------------
\section{The Mechanical Origin of Temperature}
%----------------------------------------------------------------------

\textbf{Definition:} Temperature is a direct, physical measure of the \textbf{average kinetic energy of the linear, translational motion ($v_{\text{linear}}$) of atomic vortices}:
\begin{equation}\label{eq:temp}
  T \propto \text{KE}_{\text{avg}} = \tfrac{1}{2}m_{\text{avg}}\,v_{\text{linear}}^2
\end{equation}

\textbf{Movement Budget Implication:} As temperature ($v_{\text{linear}}$) increases, the budget available for internal resonant motion ($v_{\text{orbital}}$) must decrease.

\textbf{Prediction:} At higher temperatures, $k_{\text{atom}}$ should effectively increase (as $v_{\text{orbital}}$ decreases) and ionisation energy should decrease.  This is a novel, testable prediction about the temperature dependence of atomic properties.


%----------------------------------------------------------------------
\section{The Mechanical Origin of Gas Pressure}
%----------------------------------------------------------------------

\textbf{Mechanism:} Gas pressure on container walls is the direct, mechanical result of atomic vortices \textbf{colliding with wall vortices}.

\textbf{Derivation of the Ideal Gas Law:}
\begin{itemize}
  \item Pressure $P$ = total force per area.
  \item Number of collisions/second ($N_{\text{coll}}$) $\propto$ number density ($N/V$) $\times$ average velocity ($\propto \sqrt{T}$).
  \item Average momentum change per collision ($\Delta p$) $\propto$ $m\sqrt{T}$.
  \item Combining: $P \propto (N/V) \times T$.
\end{itemize}

\begin{equation}\label{eq:ideal-gas}
  \boxed{PV = Nk_B T}
\end{equation}

Boltzmann's constant $k_B$ is revealed as a conversion factor linking vortex kinetic energy to the Kelvin scale.


%----------------------------------------------------------------------
\section{Heat and the Laws of Thermodynamics}
%----------------------------------------------------------------------

\subsection{Heat as Kinetic Transfer}

``Heat'' is the process of transferring kinetic energy ($v_{\text{linear}}$) from one ensemble of vortices to another through collisions.

\subsection{The First Law (Conservation of Energy)}

\begin{equation}
  \Delta U = Q - W
\end{equation}

A direct statement of the conservation of total movement budget in a closed system.  $U$ is the total $v_{\text{orbital}}$ budget.  $Q$ and $W$ are transfers of $v_{\text{linear}}$ budget.

\subsection{The Second Law (Entropy)}

Entropy $S$ is a measure of the \textbf{disorder and uniformity of the movement budget allocation} across a system.

\textbf{Mechanism:} In any isolated system, collisions inevitably redistribute movement budgets toward the most probable, most uniform distribution of $v_{\text{linear}}$ among all particles.  This is maximum entropy.

\textbf{The Arrow of Time:} The irreversible trend toward maximum budget uniformity.  Spontaneous re-organisation into a more ordered state would require a statistically miraculous, simultaneous, un-caused reallocation---effectively impossible.


%----------------------------------------------------------------------
\section{Phase Transitions: Catastrophic Budget Reallocation}
%----------------------------------------------------------------------

\subsection{Boiling (Liquid $\to$ Gas)}

\textbf{In liquid:} Atoms are close enough that propagated pressure fields create weak, transient ``occlusion bonds.''  $v_{\text{linear}}$ is constrained.

\textbf{At boiling point:} $v_{\text{linear}}$ budget \textbf{catastrophically overcomes} occlusion bond energy.  The system shatters.  Budget shifts almost entirely to $v_{\text{linear}}$.

\subsection{Freezing (Liquid $\to$ Solid)}

\textbf{At freezing point:} Kinetic energy drops below the binding energy of stable geometric packing.  Vortices ``lock'' into a minimum-pressure-energy crystal lattice.

\textbf{Budget shift:} $v_{\text{linear}}$ budget transfers to $v_{\text{vibrational}}$ (phonons) and is radiated away as heat.


%----------------------------------------------------------------------
\section{Conclusion}
%----------------------------------------------------------------------

The laws of thermodynamics are the emergent consequences of the \textbf{conservation and statistical distribution of the Unified Movement Budget} among vast ensembles:
\begin{itemize}
  \item \textbf{Temperature}: average linear kinetic energy of vortices.
  \item \textbf{Pressure}: physical collisions.
  \item \textbf{Entropy}: uniformity of budget allocation.
  \item \textbf{Phase transitions}: catastrophic, system-wide budget reallocations.
\end{itemize}
